% Generated semantic source for AI Almanach v3.
% Edit v3 directly after initial generation; do not overwrite
% editorial changes by rerunning the converter casually.

% ==== Scientific Writing ======================================================
\chapter{Scientific Writing and Reproducible Communication}
\chapterlead{How can a technical claim become traceable and reproducible?}{Claim}{Method}{Evidence}{Revise}{scientific-writing}

\label{ch:scientific-writing}

\section{A Paper Is a Traceable Argument}

\begin{v3topic}{Question Before Prose}
State the research question before drafting sections.
A useful question identifies the system or population, intervention or model,
comparison, outcome, and conditions under which the answer should hold.
If the question cannot be falsified or its terms cannot be measured, elegant
writing will not rescue the study.

\end{v3topic}
\begin{v3topic}{Contribution Is a Difference}
A contribution is not ``we used method X on dataset Y.''
It is the defensible difference between prior knowledge and what the study now
supports: a method, dataset, measurement, negative result, replication,
theoretical insight, or system demonstration with stated limits.

\end{v3topic}
\begin{v3topic}{Claims Have Boundaries}
Every important claim should expose its evidence, assumptions, uncertainty,
and scope.
Replace universal language with the population, dataset, time period, metric,
and operating conditions actually studied.
Separate observation (what was measured), interpretation (what it may mean),
and recommendation (what someone should do).
\end{v3topic}

\begin{figure}[htbp]
\centering
\begin{tikzpicture}[
    node distance=8mm,
    argument node/.style={draw,rounded corners=1mm,minimum width=3.2cm,
        minimum height=8mm,align=center,font=\sffamily\scriptsize},
    argument arrow/.style={-{Stealth[length=2mm]},thick}]
\node[argument node,fill=RoyalBlue!8] (claim) {Claim\\what is asserted};
\node[argument node,fill=LimeGreen!10,right=of claim] (evidence)
    {Evidence\\what was observed};
\node[argument node,fill=Orange!10,right=of evidence] (caveat)
    {Caveat\\where it may fail};
\draw[argument arrow] (evidence) -- node[above,font=\scriptsize]{supports}
    (claim);
\draw[argument arrow] (caveat) -- node[above,font=\scriptsize]{bounds}
    (evidence);
\end{tikzpicture}
\caption{A scientific argument connects a bounded claim to evidence and makes
its caveats visible.}
\label{fig:scientific-claim-evidence-caveat}
\end{figure}

\section{Design the Study Before Reporting It}

\begin{v3topic}{Protocol and Preregistration}
Write the hypotheses, primary outcomes, data exclusions, split logic, metrics,
sample-size rationale, and analysis plan before inspecting final results.
Exploratory analysis remains valuable, but label it as exploratory rather than
rewriting it as a prediction made in advance.

\end{v3topic}
\begin{v3topic}{Unit of Analysis}
Name the independent unit: person, patient, machine, site, document, session,
or time interval.
Multiple rows from one unit are not automatically independent observations.
Split and resample at the level needed to prevent leakage and
pseudoreplication.

\end{v3topic}
\begin{v3topic}{Baselines and Ablations}
Compare against a naive baseline, a strong established method, and the current
production system when one exists.
An ablation removes or changes one component to test the explanation for an
improvement.
It is not a tournament of arbitrary variants.

\end{v3topic}
\begin{v3topic}{Uncertainty and Multiplicity}
Report effect sizes and uncertainty, not only a point estimate or thresholded
$p$-value.
Account for repeated seeds, hyperparameter search, many metrics, many slices,
and repeated peeking at the test set.
The selection process is part of the experiment.
\end{v3topic}

\begin{cautionbox}[The test set is not a writing assistant]
Do not repeatedly modify the system after reading final test failures and then
report the same test set as untouched evidence.
Create a new development set, version every evaluation exposure, or collect a
fresh final set whose independence can be defended.
\end{cautionbox}

\section{Paper Architecture}

\begin{table}[htbp]
\centering
\caption{Questions answered by the major parts of an empirical AI paper.}
\label{tab:scientific-paper-architecture}
\begin{tblr}{
    width=\textwidth,
    colspec={Q[l,0.16\textwidth]X[l]X[l]},
    row{1}={font=\bfseries,bg=LimeGreen!15},
    hlines,
    vlines}
Part & Reader's question & Minimum content \\
Title and abstract & What was studied, how, and what was found? & Specific
task, method, comparison, principal result, and limitation \\
Introduction & Why is the question worth answering? & Context, gap, research
question, contributions, and scope \\
Related work & What is already known? & Organised comparison, not a citation
inventory \\
Methods & Could a competent team reproduce the study? & Data, preprocessing,
models, baselines, protocol, compute, code, and analysis \\
Results & What was observed? & Primary outcomes, uncertainty, robustness,
ablations, errors, and negative findings \\
Discussion & What do the results mean and not mean? & Interpretation,
alternative explanations, validity threats, implications, and limits \\
Artifacts & What can be inspected or rerun? & Code, data or access path,
configurations, licences, checksums, and instructions \\
\end{tblr}
\end{table}

\begin{engineernote}[Write in evidence order]
Draft the figures and tables, then the results that describe them, then the
methods needed to understand the results.
Write the introduction and abstract after the claim boundaries are stable.
This order reduces the temptation to promise a broader story than the evidence
supports.
\end{engineernote}

\section{Methods That Can Be Reproduced}

\begin{v3topic}{Data}
Report origin, collection dates, inclusion and exclusion rules, licences,
consent or legal basis, unit of analysis, label process, missingness,
preprocessing, deduplication, and split construction.
A datasheet makes assumptions and lifecycle decisions inspectable
\textsuperscript{\cite{gebruDatasheetsDatasets2021}}.

\end{v3topic}
\begin{v3topic}{Models and Training}
Name architecture, parameter count, initialization, objective, optimizer,
schedule, batch size, stopping rule, hyperparameter search, random seeds,
precision, hardware, software versions, and total compute.
Distinguish settings selected before evaluation from those selected after
examining results.

\end{v3topic}
\begin{v3topic}{Evaluation}
Define every metric mathematically or by a stable reference.
State aggregation, thresholds, slices, confidence intervals, statistical
tests, judge models or human raters, contamination checks, and how missing or
invalid outputs were handled.

\end{v3topic}
\begin{v3topic}{Artifacts and Environment}
Provide a lockfile or environment image, executable commands, configurations,
small test data, expected outputs, model and dataset identifiers, hashes, and
licences.
Archive the version used for the paper; a moving branch is not a permanent
research artifact.
\end{v3topic}

\begin{workedexample}[From inflated claim to defensible result]
\textbf{Inflated:} ``Our model significantly outperforms existing methods.''

\textbf{Defensible:} ``On the preregistered temporal test set, the proposed
model increased macro F1 from $0.712$ to $0.735$ over the tuned gradient-
boosting baseline (paired bootstrap 95\% interval for the difference:
$[0.009,0.036]$). Performance decreased below baseline at two small sites;
external prospective validation was not conducted.''

The revision names the dataset boundary, metric, baseline, effect,
uncertainty, adverse slice, and the inference the study cannot support.
\end{workedexample}

\section{Figures, Tables, and Captions as Evidence}

\begin{v3topic}{A Figure Must Answer a Question}
Choose the visual encoding from the analytical question: comparison,
distribution, relationship, change, uncertainty, or process.
Show units, denominators, sample sizes, and uncertainty.
Use consistent scales for comparisons and do not encode an important
difference with colour alone.

\end{v3topic}
\begin{v3topic}{A Caption Must Stand Alone}
A strong caption states what is shown, data and conditions, encoding, relevant
aggregation, and the conclusion the reader may draw.
Define abbreviations and mention important exclusions.
The caption should not force the reader to search several pages for the sample
size or meaning of an error bar.

\end{v3topic}
\begin{v3topic}{Tables Are for Exact Values}
Use a table when readers need exact values or multiple attributes.
Align decimals, include units in headers, state whether variation is standard
deviation, standard error, or an interval, and visually distinguish the
predeclared primary result from exploratory metrics.

\end{v3topic}
\begin{v3topic}{Accessibility Is Part of Accuracy}
Use descriptive captions or alternate text, semantic headings, readable
contrast, non-colour cues, logical reading order, and vector text where
possible.
An inaccessible result is not equally available for scrutiny.
Validate the generated PDF rather than assuming source-level semantics survive
typesetting.
\end{v3topic}

\section{Model and Dataset Reporting}

Model cards document intended use, factors, metrics, evaluation data,
limitations, and ethical considerations
\textsuperscript{\cite{mitchellModelCardsModel2019}}.
They complement rather than replace a paper: the paper argues for a result,
while a model card supports decisions about a released artifact.
Update cards when the model, data, policy, or operating environment changes.

For generative or agentic systems, also report prompt and tool interfaces,
retrieval corpus version, decoding parameters, system instructions where
disclosure is safe, evaluator configuration, refusal handling, human review,
cost and latency, safety tests, and known prompt-injection or tool-use limits.
Treat example conversations as illustrations, not as a statistically valid
evaluation.

\section{Ethics, Authorship, and AI-Assisted Writing}

\begin{v3topic}{Citation Integrity}
Read the cited source and verify that it supports the nearby claim.
Do not cite a search snippet, generated bibliography, or secondary source as
if it were the primary evidence.
Record page, section, version, and access date for unstable material.

\end{v3topic}
\begin{v3topic}{AI Assistance}
Follow the target venue's disclosure policy.
Authors remain responsible for accuracy, originality, confidentiality,
copyright, citations, and authorship criteria.
Never upload confidential manuscripts, peer-review material, personal data,
or restricted code to an unapproved service.

\end{v3topic}
\begin{v3topic}{Negative Results and Limitations}
Report credible failures, null results, excluded cases, and resource limits.
Limitations should identify how a result could break and what evidence would
be needed to extend it; ritual phrases such as ``more research is needed'' do
not help replication or decision-making.

\end{v3topic}
\begin{v3topic}{Review and Rebuttal}
Translate each review comment into a question about evidence or exposition.
Respond with a change, new analysis, bounded disagreement, or explicit
limitation.
Keep a claim--evidence matrix so edits to one result propagate to the abstract,
figures, discussion, and conclusion.
\end{v3topic}

\section{Practical Lab: From Experiment to Reproducible Report}

\begin{practicallab}[Write a compact, auditable study]
Using one experiment from an earlier chapter:
\begin{enumerate}
    \item state a falsifiable research question, primary outcome, and protocol;
    \item freeze data identifiers, splits, environment, configuration, seeds,
          and baseline before the final run;
    \item produce one figure and one table with standalone captions;
    \item report effect size, uncertainty, error slices, resource use, and one
          negative result;
    \item write a 150-word abstract containing scope, method, result, and
          limitation;
    \item create a reproducibility appendix with commands and expected output;
    \item exchange artifacts with a peer who attempts a clean rerun and records
          every missing assumption.
\end{enumerate}
Revise the claim--evidence matrix after the rerun.
Completion means another person can understand what was tested, reproduce the
main result, and see where the conclusion stops.
\end{practicallab}

\begin{checkpoint}
\begin{enumerate}
    \item What is the independent unit in your last experiment?
    \item Which sentence in your abstract has the weakest evidence?
    \item What did model selection learn indirectly from the test data?
    \item Can every figure caption be understood without the surrounding
          paragraph?
    \item Which artifact and licence would a replicator need first?
\end{enumerate}
\end{checkpoint}

%
\section{Deep Dive: Reproducibility Audit}
\begin{v3topic}{Recompute from a Clean Environment}
Give the repository and manuscript to a reader who did not create the original
environment.
Record missing data, secrets, undocumented preprocessing, manual steps,
unstable downloads, and discrepancies between figures and reported numbers.

\end{v3topic}
\begin{v3topic}{Audit the Claim Graph}
For every abstract and conclusion claim, link the supporting result, method,
artifact, and limitation.
Remove orphan claims and propagate changed numbers or caveats through captions,
tables, discussion, model cards, and supplementary material.
\end{v3topic}

\chapterreview{Claim}{Method}{Evidence}{Revise}

\section{Further Reading}

Use the venue's current author and artifact instructions as the normative
source for submission requirements.
For AI-specific documentation, model cards and datasheets provide practical
structures for reporting released models and datasets
\cite{mitchellModelCardsModel2019,gebruDatasheetsDatasets2021}.
